% ============================================================================
% Fichier  : partie4/P04_C09_cartographier.tex
% Titre    : Cartographier les Traitements
% Auteur   : MINKA MI NGUIDJOÏ Thierry Emmanuel
% Version  : 1.0 -- novembre 2026
% Corrigés : corriges/Corr_C09.tex
% ============================================================================
% Sources : loi n° 2024/017, art. 5 (fichier, interconnexion), 29, 36
% (texte officiel) ; ISO/IEC 27701 (inventaire et registre, paraphrasés).
% Faits nouveaux (bible) : 47 traitements recensés, dont 9 inconnus de la
% direction ; découverte sur les partages et les postes, 20-24 juillet 2026.
% ============================================================================

\chapter{Cartographier les Traitements}
\label{chap:cartographie}
\chaptermark{Cartographier les traitements}

\epigraph{%
    « Le registre dit ce que l'organisation croit faire.\\
    La découverte dit ce qu'elle fait. »%
}{Principe du cours}

\niveauChapitre{Fondamental}{%
    Chapitres~\ref{chap:donnees} et~\ref{chap:audit}
}

\begin{boxobjectifs}
\begin{itemize}
    \item Énoncer le contenu obligatoire du registre des activités de
          traitement (article~29).
    \item Construire un registre à partir des processus, et non des
          applications.
    \item Confronter le registre déclaré à une découverte des données
          réellement stockées.
    \item Cartographier les flux de données, y compris vers les prestataires
          et l'étranger.
    \item Mesurer la complétude d'un registre et coter ce qui reste
          inconnu.
\end{itemize}
\end{boxobjectifs}

\bigskip

% ============================================================================
\section*{Avant de commencer : la liste noire de Foumban}
% ============================================================================

Foumban, mardi 21~juillet~2026. Dans l'agence, l'équipe de KORO AUDIT
examine avec l'accord du chef d'agence le contenu du partage réseau
commun. Elle y trouve un classeur intitulé \texttt{mauvais\_payeurs.xlsx}~:
2\,140~noms, numéros de téléphone, quartiers, montants impayés, et une
colonne « remarques » où figurent des mentions comme « famille du chef de
quartier, ne pas relancer » ou « commerçant bamoun, fiable malgré
retard ». Le fichier est partagé entre six agences par messagerie. Il ne
figure dans aucune liste de l'établissement.

Ce classeur est un traitement à part entière~: il a une finalité (décider à
qui prêter), des personnes concernées, des destinataires. Il contient des
données sensibles (une origine régionale et des informations bancaires), et
il sert à décider sans que les personnes le sachent.

\begin{boxalert}
\textbf{Violation typique.} Établir le registre par questionnaire auprès des
directions, puis le tenir pour complet.

\smallskip
\textbf{Règle~:} un registre déclaratif se confronte à une découverte des
données réellement stockées et des flux réellement observés. Ce que la
découverte trouve et que le registre ignore mesure l'incomplétude du
registre.
\end{boxalert}

% ============================================================================
\section{Le Registre de l'Article 29}
\label{sec:c04-registre}
% ============================================================================

\begin{boxjuris}
\textbf{\loicam{n°~2024/017}, article~29}~\cite{loi2024017}~:
\textit{« (1) Le responsable du traitement ou le sous-traitant tient un
registre physique ou numérique des activités de traitement effectuées sous
leur responsabilité. (2) Le registre visé à l'alinéa 1 ci-dessus comporte
les informations ci-après~: le nom et les coordonnées du responsable du
traitement et, le cas échéant, le nom du sous-traitant~; les finalités du
traitement~; la description des catégories de personnes concernées et des
catégories des données à caractère personnel~; les catégories de
destinataires auxquels les données à caractère personnel ont été ou seront
communiquées~; les documents attestant de l'existence de garanties
appropriées ou le numéro de l'autorisation de l'Autorité de protection des
données à caractère personnel. »}

\smallskip
\textbf{Portée.} L'obligation pèse sur le responsable \textit{et} sur le
sous-traitant. Le registre peut être papier ou numérique. Son dernier champ
relie le registre aux autorisations~: un traitement sans numéro
d'autorisation ni garanties documentées est, par construction, un écart.
\end{boxjuris}

Le contenu légal est un minimum. Un registre utile à l'audit et à la gestion
ajoute la base du traitement (consentement ou dérogation de l'article~9), la
durée de conservation, les transferts hors du Cameroun, les mesures de
sécurité, la date de la dernière revue et le propriétaire du traitement dans
l'organisation.

\begin{tableaucours}{L{4.0cm}cL{6.6cm}}{%
    \tableauHeader{Champ} &
    \tableauHeader{Obligatoire} &
    \tableauHeader{Ce que l'auditeur vérifie}
}
Responsable, sous-traitant & Oui & Coordonnées exactes~; sous-traitants
cohérents avec les contrats \\
Finalités & Oui & Une finalité par traitement, déterminée et explicite \\
Personnes et données & Oui & Catégories complètes, données sensibles
signalées \\
Destinataires & Oui & Comparés aux flux observés \\
Garanties ou autorisation & Oui & Numéro d'autorisation ou pièces
justificatives \\
Fondement & Recommandé & Cohérent avec l'article~9 \\
Durée de conservation & Recommandé & Cohérente avec l'article~13 \\
Transferts & Recommandé & Pays, garanties (Chapitre~\ref{chap:transferts}) \\
Mesures de sécurité & Recommandé & Renvoi aux mesures de l'article~27 \\
Propriétaire, date de revue & Recommandé & Une personne, une date récente \\
\end{tableaucours}
\vspace{-0.8em}
\begin{center}{\small\itshape Tableau~9.1. Les champs du registre}\end{center}

% ============================================================================
\section{Recenser par les Processus}
\label{sec:c04-processus}
% ============================================================================

L'erreur la plus courante consiste à partir de l'inventaire des
applications. Une application sert souvent plusieurs traitements (le
logiciel bancaire sert la tenue des comptes, le crédit, le recouvrement et
la lutte contre le blanchiment)~; à l'inverse, un traitement peut ne
reposer sur aucune application (un cahier, un classeur partagé, un groupe de
messagerie).

On part donc des \termecle{processus} de l'organisation~: pour chaque
direction, quelles activités, quelles personnes concernées, quelles
décisions prises sur ces personnes. Le Chapitre~\ref{chap:donnees} a fourni
les familles de données~; on les croise ici avec les processus pour obtenir
les traitements, puis on rattache à chacun les applications et les supports
qui le portent.

\begin{boxprocedure}
\textbf{Établir le registre d'une organisation.}
\begin{enumerate}
    \item Lister les processus de chaque direction, avec un interlocuteur.
    \item Pour chaque processus, identifier les personnes concernées et les
          décisions prises sur elles.
    \item En déduire les traitements, un par finalité.
    \item Rattacher à chaque traitement ses supports~: applications,
          fichiers bureautiques, papier, messagerie.
    \item Renseigner les champs du registre, avec la pièce qui fonde chaque
          champ.
    \item Confronter le registre à la découverte et aux flux
          (sections~\ref{sec:c04-decouverte} et~\ref{sec:c04-flux}).
    \item Faire valider chaque fiche par son propriétaire, et dater la
          revue.
\end{enumerate}
\end{boxprocedure}

% ============================================================================
\section{La Découverte des Données}
\label{sec:c04-decouverte}
% ============================================================================

La \termecle{découverte des données} consiste à chercher, dans les supports
réels, des données personnelles que le registre ne mentionne pas. Elle se
pratique sur quatre types de supports.

\begin{itemize}
    \item \textbf{Les partages de fichiers et les postes}~: recherche de
          motifs (numéros de téléphone, numéros de compte, numéros de carte
          nationale d'identité) dans les documents bureautiques.
    \item \textbf{Les bases de données}~: lecture des schémas, à la
          recherche de colonnes nommées \texttt{nom}, \texttt{tel},
          \texttt{cni}, \texttt{date\_naissance}.
    \item \textbf{La messagerie et les outils de collaboration}~: groupes,
          pièces jointes récurrentes, formulaires en ligne.
    \item \textbf{Le papier}~: armoires, cahiers, archives d'agence.
\end{itemize}

La découverte est elle-même un traitement de données personnelles. Elle se
fait avec l'accord écrit du commanditaire, sur les seuls supports de
l'organisation, en ne relevant que des indicateurs (nom du fichier,
emplacement, nombre d'occurrences, types de données) et jamais le contenu
lui-même, sauf échantillon nécessaire à la qualification.

\begin{boxinfo}
\textbf{Les faux positifs.} Une recherche de motifs trouve des numéros qui
n'en sont pas~: un numéro de facture peut ressembler à un numéro de
téléphone. On qualifie donc chaque fichier signalé par un examen rapide d'un
échantillon, et l'on cote la proposition « ce fichier contient des données
personnelles » \OUI{}, \NON{} ou \NEK{}. Le laboratoire APDC-LAB-02 exerce
cette qualification.
\end{boxinfo}

% ============================================================================
\section{Cartographier les Flux}
\label{sec:c04-flux}
% ============================================================================

Un registre décrit des traitements~; une \termecle{cartographie des flux}
décrit le chemin des données entre les traitements, les systèmes, les
prestataires et les pays. Elle répond à trois questions que le registre
laisse ouvertes~: par où les données entrent-elles, où sont-elles stockées,
et vers qui sortent-elles.

On l'établit à partir de trois sources indépendantes~: les entretiens
(ce que les équipes décrivent), les configurations (connecteurs,
exportations planifiées, interfaces de programmation) et les traces
réseau ou les factures (ce qui sort réellement, et vers qui l'organisation
paie un service). Une facture mensuelle d'un service de formulaires en
ligne révèle souvent un flux que personne n'avait déclaré.

L'article~36 encadre l'\termecle{interconnexion des fichiers}, définie à
l'article~5 comme la mise en relation de données traitées pour une
finalité avec d'autres données. Elle n'est licite que si elle est
« nécessaire au respect d'une obligation légale », et ne doit entraîner ni
discrimination ni réduction des droits. Le classeur des mauvais payeurs,
alimenté par six agences et rapproché du fichier des clients, est une
interconnexion~; l'auditeur en cherche l'obligation légale, et n'en trouve
pas.

% ============================================================================
\section{Le Registre de NYANGA}
\label{sec:c04-nyanga}
% ============================================================================

\begin{boxscenario}{Mission NYANGA. Quarante-sept traitements, 20 au 24 juillet 2026}
\textbf{Constatations.}
\begin{itemize}
    \item Les entretiens avec les huit directions produisent un registre
          déclaratif de 38~traitements.
    \item La découverte sur les partages du siège et des trois agences signale
          1\,312 fichiers contenant des motifs de données personnelles~;
          l'examen d'un échantillon de 120 fichiers en confirme 97.
    \item Neuf traitements apparaissent que personne n'a déclarés~: le
          classeur des mauvais payeurs~; deux groupes de messagerie où les
          agents échangent des photos de pièces d'identité~; un questionnaire
          de satisfaction hébergé par un service de formulaires étranger~; un
          kit de mesure d'audience intégré à l'application mobile, qui envoie
          les identifiants des téléphones à un tiers~; l'enregistrement des
          appels du service client~; le fichier des clients « privilégiés »
          du service marketing~; les badges d'accès du siège~; les CV reçus
          par messagerie.
    \item Les factures de 2025 révèlent trois prestataires absents de la
          liste des sous-traitants.
\end{itemize}

\textbf{Analyse} (confiance élevée). Le registre complet compte
47~traitements. Le registre déclaratif en couvrait 38, soit
81~\%~; les neuf manquants sont ceux qui transfèrent des données à des
tiers ou qui prennent des décisions sans procédure, c'est-à-dire les plus
risqués.
\end{boxscenario}

\begin{boxobjet}{Liste des mauvais payeurs, agences de l'Ouest}
\textbf{Finalité réelle~:} décider de l'octroi de crédit et de la relance.
\textbf{Personnes~:} 2\,140 débiteurs, dont des clients d'autres
établissements cités par des tiers. \textbf{Données~:} identité, téléphone,
quartier, montant impayé, remarques incluant l'origine régionale (donnée
sensible). \textbf{Destinataires~:} six agences, par messagerie.
\textbf{Ce que l'audit établit~:} aucun fondement (\NON{}), interconnexion
sans obligation légale (\NON{}), aucune information des personnes
(\NON{})~; usage pour des décisions de crédit (\OUI{}, entretiens et
échantillon de 12 refus).
\end{boxobjet}

\begin{boxlecture}{Le registre de NYANGA, 24 juillet 2026}
\begin{tableaucours}{L{7.8cm}cL{3.2cm}}{%
    \tableauHeader{Exigence} &
    \tableauHeader{Cote} &
    \tableauHeader{Preuve}
}
Un registre existe & \OUI{} & Établi par l'audit \\
Le registre est complet & \NON{} & 9 traitements découverts \\
Chaque fiche comporte les champs de l'article~29 & \NON{} & Garanties
absentes \\
Chaque traitement a un propriétaire & \NON{} & 14 sans propriétaire \\
Les interconnexions sont fondées sur une obligation légale & \NON{} &
Classeur des mauvais payeurs \\
Les flux vers les tiers sont tous connus & \NEK{} & Trois prestataires
découverts, recherche inachevée \\
Les fichiers papier des agences non visitées sont recensés & \NEK{} & 19
agences hors échantillon \\
\end{tableaucours}

\medskip
\textbf{Décompte.} Un \OUI{}, quatre \NON{}, deux \NEK{} sur sept.
\hfill\lecture{0{,}14}{0{,}29}{1{,}38}

\medskip
\textbf{Lecture.} Le registre existe désormais, mais c'est l'audit qui l'a
établi~: le \OUI{} ne crédite pas l'établissement. Avec $U = 0{,}29$, la
complétude reste opaque~; l'incertitude se localise dans les flux et dans
les dix-neuf agences non visitées.
\end{boxlecture}

% ============================================================================
\section{Ce que la Cartographie demande à chaque Métier}
\label{sec:c04-metiers}
% ============================================================================

\begin{boxmetier}{Auditeur de la protection des données}
Mesurer la complétude du registre par la découverte et les flux, et
rapporter le taux de couverture. Un registre à 81~\% n'est pas un registre
« presque complet »~: les traitements manquants sont rarement les moins
risqués.
\end{boxmetier}

\begin{boxmetier}{Délégué à la protection des données}
Tenir le registre comme un outil vivant~: une fiche par traitement, un
propriétaire, une date de revue, et une procédure qui oblige tout nouveau
projet à créer sa fiche avant de collecter.
\end{boxmetier}

\begin{boxmetier}{RSSI}
Fournir les sources de la découverte~: exports des schémas de bases,
journaux de flux sortants, inventaire des services en ligne payés. Bloquer
les outils de partage qui contournent les traitements déclarés.
\end{boxmetier}

\begin{boxresume}
\begin{itemize}
    \item L'article~29 impose un registre au responsable et au
          sous-traitant, avec cinq catégories d'informations au minimum.
    \item On recense par les processus et les personnes concernées, puis on
          rattache les applications et les supports.
    \item Le registre déclaratif se confronte à la découverte des données et
          aux flux réels~; l'écart mesure son incomplétude.
    \item L'interconnexion de fichiers n'est licite que si une obligation
          légale l'impose (article~36).
    \item À NYANGA, 47~traitements, dont 9 inconnus de la direction et parmi
          eux un fichier de mauvais payeurs mentionnant l'origine régionale.
\end{itemize}
\end{boxresume}

\begin{boxinfo}
\textbf{Points de contrôle.} Vous avez compris ce chapitre si vous savez~:
\begin{itemize}
    \item citer les champs obligatoires du registre~;
    \item expliquer pourquoi un inventaire d'applications n'est pas un
          registre~;
    \item mesurer la complétude d'un registre à partir d'une découverte.
\end{itemize}
\end{boxinfo}

\begin{boxexo}[title={\ding{45}\quad Exercice~9.1. Traitement ou application \niveaubase}]
Dans la liste suivante, distinguez les traitements des moyens
techniques~: (a) la paie~; (b) le serveur de messagerie~; (c) le
recrutement~; (d) la base Oracle~; (e) la vidéosurveillance des agences~;
(f) le logiciel bancaire.
\end{boxexo}

\begin{boxexo}[title={\ding{45}\quad Exercice~9.2. Une fiche de registre \niveaumoyen}]
Rédigez la fiche de registre du traitement « enregistrement des appels du
service client », avec les champs obligatoires et recommandés. Signalez les
champs que vous ne pouvez pas remplir sans pièce supplémentaire.
\end{boxexo}

\begin{boxexo}[title={\ding{45}\quad Exercice~9.3. Mesurer la complétude \niveaumoyen}]
La découverte d'une autre organisation signale 800 fichiers, dont un
échantillon de 50 donne 35 confirmations. Le registre déclaratif compte
20~traitements~; la découverte en ajoute 6. (a) Quel est le taux de
couverture du registre déclaratif~? (b) Combien de fichiers contiennent
probablement des données personnelles, et avec quelle réserve~?
\end{boxexo}

\begin{boxexo}[title={\ding{45}\quad Exercice~9.4. La liste noire \niveauexpert}]
Le directeur des crédits défend le classeur des mauvais payeurs~: « Sans
lui, nous prêterions à des fraudeurs. » Analysez sa licéité au regard des
articles~5, 9, 11, 17, 36 et~48. Proposez une solution qui réponde au
besoin de l'établissement dans le respect de la loi.
\end{boxexo}

\bigskip

\begin{boxinfo}
\textbf{Pour aller plus loin.}
\begin{itemize}
    \item L'ISO/IEC~27701~\cite{iso27701}, pour la tenue des inventaires et
          des registres dans un système de management de la vie privée.
\end{itemize}
\end{boxinfo}

\begin{boxinfo}
\textbf{Transition.} Quarante-sept traitements ne présentent pas le même
danger pour les personnes. Lesquels examiner d'abord, et lesquels exigent
une évaluation d'impact~? C'est l'objet du Chapitre~\ref{chap:risque}.
\end{boxinfo}
